% AUTO-GENERATED by papers/tables/generate_tables_ch4_overview.py - do not edit by hand.
\begin{table}[htbp]
  \centering
  \begin{footnotesize}
  \caption{Extended overview summary with novelty-kind and similarity diagnostics.}
  \label{tab:overview_extended_summary}
  \begin{threeparttable}
  \begin{tabularx}{\textwidth}{@{}L{0.70} S[table-format=2.2] S[table-format=2.2] S[table-format=1.3] S[table-format=1.3] L{1.30}@{}}
    \toprule
    \textbf{Test} & {\textbf{Novel (\%)}} & {\textbf{Novel-kind (\%)}} & {\textbf{Mean Sim.}} & {\textbf{SD Sim.}} & \textbf{Traceability metric} \\
    \midrule
    T1 & 0.00 & 5.90 & 0.305 & 0.078 & traceability flag \\
    T2 & 4.76 & 4.76 & 0.581 & 0.072 & training-derived flag \\
    T3 & 27.38 & 32.03 & 0.672 & 0.051 & likely-derivative flag \\
    T4 & 63.69 & 59.13 & 0.462 & 0.031 & training-derived flag \\
    \bottomrule
  \end{tabularx}
  \begin{tablenotes}
    \item Novel-kind retains each test's secondary signal (continuous where available; binary proxy otherwise).
  \end{tablenotes}
  \end{threeparttable}
  \end{footnotesize}
\end{table}
